\chapter{PRISTUP TEMELJEN NA UČENJU}
\label{ch:rl}

Pristup temeljen na učenju rješava isti zadatak kao poglavlje \ref{ch:klasicni}, ali bitno
drugačijim sredstvima. Klasični pristup vrh alata vodi krutim pozicijskim
upravljanjem po unaprijed izračunatoj putanji, uz geometriju ograničenja poznatu iz percepcije.
Ovdje se, umjesto toga, uvodi varijabilna impedancija: politika u svakom koraku sama bira koliko
će regulator biti krut po svakoj osi, a geometriju ograničenja ne dobiva ni iz percepcije ni iz
eksplicitnog procjenitelja, nego je mora otkriti iz vlastitih opažanja sile i gibanja. To je
izravan odgovor na nalaz iz poglavlja \ref{ch:klasicni}: kruto pozicijsko praćenje luka davalo je
velike kontaktne sile na svako odstupanje procjene, pa varijabilna impedancija postoji upravo
zato da politika tu krutost regulira sama.

\section{Opseg}
\label{sec:rl-opseg}

Politika ne uči prilaz ni hvat kvake. Epizoda počinje od stanja u kojem prihvatnica već drži
kvaku, a klasični postupak prilaza i hvata (potpoglavlje \ref{sec:prilaz}) smatra se zajedničkim
za oba pristupa. Razlog je dvostruk: prilaz i hvat rješavaju problem navigacije i planiranja
gibanja, a ne problem interakcije koji je predmet ovog rada, i njihovo bi dupliciranje otežalo
poštenu usporedbu, jer bi mjerena razlika između pristupa uključivala i razliku u tome tko je
bolji u navigaciji, ne samo u interakciji. Iz istog je razloga izostavljen i prolazak kroz vrata:
i on je navigacija kroz poznatu geometriju, ista kategorija kao prilaz. Ta odluka o opsegu izravno
oblikuje ono što se u poglavlju \ref{ch:rezultati} smije uspoređivati. Usporedba se odnosi
isključivo na fazu otvaranja, ne na cijeli zadatak, dok klasični pristup taj cijeli zadatak
stvarno izvodi od početka do kraja.

Sam postupak prolaska kroz vrata je implementiran i radi, uključivo za slučaj naučene politike,
ali je namjerno isključen iz rezultata radi dosljednosti opsega. Vrijedan je kao potvrda da se
naučena politika može ugraditi u isti izvršni sloj kao klasični pristup.

Trening se u cijelosti odvija unutar simulacijskog okruženja Isaac Lab, bez sustava ROS 2 u
regulacijskoj petlji. Razlog je paralelizacija: trening istovremeno pokreće više stotina
okruženja, a posredovanje kroz ROS 2 bilo bi usko grlo. Epizoda traje 10 sekundi, uz korak
upravljanja od 60 Hz (fizika 120 Hz, decimacija 2), što daje 600 koraka po epizodi.

Uspješnost politike pri determinističkom izvođenju, koje se razlikuje od učenja, mjerena je
zasebno i uspoređena s klasičnim pristupom u poglavlju \ref{ch:rezultati}.

\section{Formulacija Markovljevog procesa odlučivanja}
\label{sec:mdp}

\subsection{Prostor opažanja}
\label{sec:opazanje}

Vrijedi isto načelo kao u klasičnom pristupu: kut ili pomak zgloba vrata, dakle stvarno stanje
zadatka, nikada nije dio opažanja politike, nego se koristi isključivo u nagradi i terminaciji,
gdje je legitiman jer se računa izvan regulacijske petlje. Opažanje sadrži:

\begin{itemize}
    \item pozu vrha alata u koordinatnom sustavu platforme (položaj i orijentacija);
    \item linearnu i kutnu brzinu vrha alata, u istom koordinatnom sustavu;
    \item brzinu platforme, (linearna brzina po $x$ i $y$, kutna brzina oko $z$);
    \item procijenjenu silu i moment na vrhu alata, uz Gaussov šum standardne devijacije
          0.5 N, koji odgovara šumu stvarnog procjenitelja iz poglavlja \ref{ch:procjena};
    \item vektore od platforme do najbližih točaka susjednih prepreka (zid, dovratnik), uz
          Gaussov šum standardne devijacije 1 cm, koji odgovara izmjerenoj stabilnosti sredine
          otvora iz skena lidara u klasičnom pristupu (potpoglavlje \ref{sec:prolazak});
    \item posljednju akciju.
\end{itemize}

\subsection{Prostor akcija}
\label{sec:akcije}

Ruka se upravlja regulatorom u operacijskom prostoru s varijabilnom impedancijom. Akcija zadaje
pomak reference vrha alata u odnosu na trenutnu pozu, za svih šest stupnjeva slobode, i šest
pripadnih krutosti, po jedna za svaku os. Time politika izravno bira, u svakom koraku i za
svaku os zasebno, koliko će regulator biti krut, točno ono što je poglavlje
\ref{ch:klasicni} pokazalo da kruto pozicijsko upravljanje ne može samo po sebi ostvariti.

Sam pomak reference po koraku ograničen je na 15 mm. Pri većoj vrijednosti, istraživačko ponašanje
na početku učenja trgalo je hvat sa šipke prije nego što bi politika uopće stigla nešto naučiti.

Uz šest stupnjeva slobode ruke, akcija sadrži i tri komponente brzine platforme (uzdužnu,
poprečnu i kutnu) pa politika sama koordinira ruku i platformu.

\subsection{Koordinacija platforme i manipulatora}
\label{sec:baza-uci}

Uvedena je dodatna kazna na odstupanja platforme od okomice na
ravninu vrata, s težinom razmjernom tome koliko takvo odstupanje kod danog tipa vrata stvarno
šteti: znatno većom kod zakretnih, gdje dijagonalan položaj platforme kasnije ne prolazi kroz
suženi otvor, nego kod kliznih, gdje platforma svejedno mora pratiti krilo poprečno pa okomitost
manje košta. Kazna je uz to zatvorena ovisno o ostvarenom napretku, tako da ne djeluje u ranoj
fazi zadatka, kada bi se sukobila s potrebom da politika zakrene platformu radi dosega. Ruka
je montirana izvan osi platforme, pa zakret platforme proširuje radni prostor ruke na legitiman način.

\subsection{Funkcija nagrade}
\label{sec:nagrada}

Tablica \ref{tab:nagrada} navodi sve članove nagrade. Sve su težine kalibrirane mjerenjem, ne
pogođene: cilj je bio da svaki član, u konvergiranom pokretanju, iznosi između 3 i 10 \% udjela
glavnog člana napretka, a postupak je bio pustiti probno pokretanje, očitati udjele i skalirati
težinu razmjerno odstupanju od tog cilja.

\begin{table}[H]
    \centering
    \caption{Članovi funkcije nagrade}
    \label{tab:nagrada}
    \begin{tabular}{@{}llp{0.5\textwidth}@{}}
        \toprule
        Član                 & Težina                                 & Svrha                                                   \\
        \midrule
        napredak             & $+10{,}0$                              & glavni član, normiran na prag uspjeha                   \\
        okomita sila         & $-0{,}02$                              & kazni silu okomitu na stvarni smjer gibanja vrata       \\
        zasićenje zglobova   & $-2{,}0$ (klizna), $-6{,}0$ (zakretna) & kazni rad blizu granice
        raspona zglobova ruke                                                                                                   \\
        singularitet         & $-20{,}0$                              & kazni pad manipulabilnosti ispod praga                  \\
        ubrzanje ruke        & $-5\times 10^{-4}$                     & kazni ubrzanje zglobova ruke                            \\
        trzaj ruke           & $-1\times 10^{-7}$                     & kazni trzaj (drugu derivaciju brzine) zglobova ruke     \\
        prekomjerna sila     & $-0{,}05$                              & kazni ukupnu silu iznad referentne granice              \\
        poravnanje hvata     & $-3{,}0$                               & kazni odstupanje prilazne osi hvata od okomice na krilo \\
        poravnanje platforme & $-6{,}0$ (zakretna), $-0{,}1$ (klizna) & kazni odstupanje
        platforme od okomice na ravninu vrata                                                                                   \\
        prodor u krilo       & $-20{,}0$                              & kazni približavanje ruba platforme krilu ispod
        sigurnog razmaka                                                                                                        \\
        prodor u zid         & $-20{,}0$                              & kazni približavanje ruba platforme zidu ili dovratniku
        ispod sigurnog razmaka                                                                                                  \\
        promjena akcije      & $-0{,}002$                             & kazni razliku uzastopnih akcija                         \\
        uspjeh               & $+2{,}0$                               & bonus po koraku dok je pomak vrata iznad praga uspjeha  \\
        \bottomrule
    \end{tabular}
\end{table}

Napredak se ne normira na puni hod zgloba vrata, već na puno manju referentnu vrijednost
0.6 m kod kliznih i 0.5 rad (28.6 \textdegree) kod zakretnih vrata. Uz normiranje na puni hod,
manji pomak nosio bi premalu nagradu da nadjača kazne, pa bi politika naučila dovesti vrata na
određeni pomak i ne pomicati ih dalje. Time napredak nije ograničen odozgo, već smije prijeći jedinicu,
što znači da nagrada za napredak raste što su vrata više otvorena.

Prag uspjeha, koji određuje kada se dodjeljuje bonus iz zadnjeg retka tablice \ref{tab:nagrada},
zasebna je veličina. Kod kliznih vrata poklapa se s referentnom vrijednošću normalizacije i iznosi
0.6 m. Kod zakretnih vrata je namjerno odvojen i postavljen na 1.57 rad (90 \textdegree).
Razlog je taj što je 0.5 rad dovoljno da se vrata mjerljivo pomaknu, ali nije dovoljno da
robot kroz njih prođe. Normalizacija je pritom ostala nepromijenjena kako bi kalibracija
svih ostalih veličina ostala valjana. Bonus za uspjeh dodjeljuje se po koraku dok je pomak
iznad praga, a ne kao jednokratna nagrada uz prekid epizode.

Kazna na singularitet koristi Yoshikawin indeks manipulabilnosti \cite{yoshikawa1985manipulability}:

\begin{equation}
    w(\bm{q}) = \sqrt{\det\left( \bm{J}(\bm{q}) \bm{J}(\bm{q})\transp \right)} ,
    \label{eq:manip}
\end{equation}

koji je jednak nuli točno na kinematičkom singularitetu i nigdje drugdje. Kazna je normirana na
raspon od nule do jedan.

\subsection{Terminacije}
\label{sec:terminacije}

Epizoda prekida prije isteka vremena u četiri slučaja

\begin{itemize}
    \item sila je prešla sigurnosnu granicu,
    \item hvat je izgubljen jer se vrh alata previše udaljio od šipke kvake,
    \item rub platforme je dotaknuo krilo vrata,
    \item rub platforme dotaknuo zid ili dovratnik.
\end{itemize}

\subsection{Nasumičnost pri resetu i kurikulum otpora}
\label{sec:randomizacija}

Pri svakom resetu, kut prilaska platforme vrtnji se nasumično zakreće unutar $\pm$ 20\textdegree{}
oko nazivne vrijednosti. Otpor vrata randomizira se u rasponima navedenima u poglavlju
\ref{ch:sustav} (tablica \ref{tab:randomizacija}).

Kod kliznih vrata otpor se uz to uvodi postupno, kurikulumom koji tijekom prvih 1500 iteracija
učenja linearno širi gornju granicu trenja i prigušenja s vrijednosti koja odgovara laganim
vratima do pune vrijednosti iz tablice \ref{tab:randomizacija}, nakon čega ostaje na punoj
vrijednosti do kraja treninga. Kurikulum politici omogućuje da najprije nauči na lakšim vratima, gdje
se povratna informacija o napretku redovito javlja, prije nego što se suoči s punim rasponom.
Zakretna vrata kurikulum nemaju, jer je njihov otpor za red veličine manji i politika na njemu
uči bez poteškoća; različit režim treninga po tipu vrata posljedica je te razlike u težini
zadatka, ne proizvoljne odluke.

\section{Postupak učenja}
\label{sec:trening}

Učenje se provodi algoritmom proksimalne optimizacije politike (PPO) \cite{schulman2017ppo}, s
hiperparametrima u tablici \ref{tab:ppo}.

PPO pripada familiji postupaka s aktorom i kritikom, koji uče dvije odvojene neuronske mreže.
Aktor je sama politika, ona iz opažanja daje srednje vrijednosti normalne razdiobe iz koje se potom
uzorkuje akcija, po jednu za svaku dimenziju akcijskog prostora. Kritik ničim ne upravlja, nego
iz opažanja procjenjuje očekivani ukupni diskontirani povrat iz
zatečenog stanja, dakle mjeru toga koliko je stanje samo po sebi povoljno. Kritik je potreban
jer ostvarena nagrada sama po sebi ne govori je li akcija bila dobra: visoka nagrada može biti
posljedica povoljnog stanja, a ne odabrane akcije. Razlika između stvarno ostvarenog povrata i
kritikove procjene, tzv. prednost, ono je što ulazi u ažuriranje aktora. Pozitivna
prednost pojačava vjerojatnost odabrane akcije, negativna je smanjuje. Diskontni faktor $\gamma$
i faktor generalizirane procjene prednosti $\lambda$ iz tablice \ref{tab:ppo} upravljaju upravo
izračunom te prednosti.

Obje su mreže potpuno povezane i međusobno neovisne, bez zajedničkih slojeva, s trima skrivenim
slojevima širine 256, 128 i 64. Aktor preslikava 41 ulaznu veličinu, koliko ih ima opažanje iz
potpoglavlja \ref{sec:opazanje}, u 15 izlaznih, koliko ih ima akcija iz potpoglavlja
\ref{sec:akcije}. Kritik ima isti ulaz, ali jedan jedini izlaz. Postupno suženje slojeva
uobičajen je obrazac: široki prvi sloj izdvaja značajke iz sirovog opažanja, a svaki ih sljedeći
sažima prema odluci. Ukupno je riječ o svega stotinjak tisuća parametara po obje mreže, što je
za mjerila neuronskih mreža malo, ali uobičajeno kada je opažanje vektor izmjerenih veličina, a
ne primjerice slika.

Kao aktivacijska funkcija skrivenih slojeva koristi se ELU (engl. \textit{exponential linear
    unit}),

\begin{equation}
    \mathrm{ELU}(x) =
    \begin{cases}
        x,         & x > 0    \\
        e^{x} - 1, & x \leq 0
    \end{cases}
    \label{eq:elu}
\end{equation}

koja unosi nelinearnost, bez koje bi tri uzastopna linearna sloja bila jednakovrijedna jednomu.
Pozitivne ulaze ne mijenja, dok za negativne daje glatku negativnu
vrijednost umjesto nule. Time se, prvo, izbjegava da neuron s negativnim izlazom ostane trajno
neaktivan, jer mu gradijent ne pada na nulu, i drugo, srednja vrijednost aktivacija ostaje bliža
nuli, što ubrzava konvergenciju. Potonje je ovdje osobito prikladno, jer su ulazne veličine
(sile, brzine i pomaci) predznačene i prirodno raspoređene oko nule, pa bi funkcija koja
negativne vrijednosti postavlja na nulu odbacila polovicu njihova raspona.

\begin{table}[H]
    \centering
    \caption{Hiperparametri učenja}
    \label{tab:ppo}
    \begin{tabular}{@{}ll@{}}
        \toprule
        Hiperparametar                                     & Vrijednost                                   \\
        \midrule
        Broj koraka po okruženju prije ažuriranja          & 24                                           \\
        Broj epoha učenja po ažuriranju                    & 5                                            \\
        Broj mini-grupa                                    & 4                                            \\
        Stopa učenja                                       & $10^{-3}$, prilagodljiva prema odstupanju KL \\
        Diskontni faktor $\gamma$                          & 0.99                                         \\
        Faktor generalizirane procjene prednosti $\lambda$ & 0.95                                         \\
        Koeficijent obrezivanja                            & 0.2                                          \\
        Koeficijent entropije                              & 0.001                                        \\
        Početna standardna devijacija politike             & 0.5                                          \\
        \bottomrule
    \end{tabular}
\end{table}

Broj paralelnih okruženja postavljen je na 1024, kao kompromis između propusnosti i broja
dostupnih pokretanja. Svako je pokretanje trajalo 2400 iteracija, a sve brojke u nastavku očitane su kao
prosjek zadnjih 100 iteracija, zato što pojedinačna iteracija nosi previše šuma da bi sama
predstavljala izvedbu politike.

\section{Rezultati}
\label{sec:rl-rezultati}

Tablica \ref{tab:rl-rezultati} navodi rezultate na sceni sa zidovima i s graničnim
vrijednostima zgloba vrata usklađenima s klasičnim pristupom. Radi usporedivosti, obje su
politike trenirane s istim brojem okruženja i istim brojem iteracija, a sve su vrijednosti
prosjek zadnjih 100 iteracija.

Prva dva retka tablice vrijednosti su pripadnih članova nagrade, a ne veličine koje izravno
opisuju ishod. Okruženje ih prije zapisivanja dijeli duljinom epizode, pa su ispod svakog člana
navedene veličine izvedene iz njega: član napretka razmjeran je prosječnom otvoru vrata kroz
epizodu, a član uspjeha udjelu epizode tijekom kojeg su vrata iznad praga uspjeha. Udjeli
terminacija u nastavku tablice jesu udjeli epizoda i zbrajaju se na jedinicu.

\begin{table}[H]
    \centering
    \caption{Rezultati pristupa temeljenog na učenju, 1024 okruženja, 2400 iteracija}
    \label{tab:rl-rezultati}
    \begin{tabular}{@{}lcc@{}}
        \toprule
        Veličina                           & Zakretna vrata & Klizna vrata \\
        \midrule
        Član napretka                      & 18.93          & 9.87         \\
        \quad prosječni otvor kroz epizodu & 0.947 rad      & 0.592 m      \\
        Član uspjeha                       & 0.743          & 0.831        \\
        \quad udio epizode iznad praga     & 0.371          & 0.416        \\
        \midrule
        Udio: izgubljen hvat               & 0.072          & 0.000        \\
        Udio: dodir krila                  & 0.061          & 0.003        \\
        Udio: dodir zida                   & 0.002          & 0.000        \\
        Udio: prekomjerna sila             & 0.006          & 0.028        \\
        Udio: istek vremena                & 0.860          & 0.969        \\
        Duljina epizode, koraci            & 561.8          & 587.4        \\
        \bottomrule
    \end{tabular}
\end{table}

Te vrijednosti opisuju ponašanje politike tijekom učenja, uz šum istraživanja. Stopa uspjeha pri
determinističkom izvođenju mjerena je zasebno i navedena je u poglavlju \ref{ch:rezultati}.

\section{Ograničenja}
\label{sec:rl-limits}

\begin{itemize}
    \item Naučena politika koristi srednju vrijednost razdiobe akcije, bez šuma istraživanja, pa
          se izvedba pri izvođenju ne mora poklapati s brojkama iz treninga, koje su mjerene uz
          šum.

    \item Kazna za odstupanje platforme od okomice na vrata pogoršava se u zadnjoj četvrtini
          treninga, jer okomit položaj platforme stvarno smanjuje doseg ruke montirane izvan
          njezine osi, pa politika pred kraj treninga taj kompromis rješava nešto drugačije nego
          ranije.

    \item Politika, za razliku od klasičnog pristupa, ne raspolaže nikakvim posredovanim
          signalom o tipu ili smjeru ograničenja: geometriju mora otkriti iz iste sile i pomaka
          kojima i djeluje.
\end{itemize}